\chapter{结论与展望}

\section{主要研究结论}

本研究围绕深远海可再生能源的系统集成与多元利用，完成了产业属性分析、系统架构设计、
源储算用协同调度建模、典型时段与年度仿真，以及投资和政策机制研究。主要结论如下。

\begin{enumerate}[label=(\arabic*)]
  \item 深远海能源岛具有未来产业的主要特征。漂浮式风电、构网型储能、海上制氢、绿色算力
        和智能调度等技术已经具备不同程度的示范基础，但完整系统仍处于集成验证阶段。其规模化
        发展依赖深远海环境下的长期可靠性、跨专业接口、产品合同和项目融资条件。
  \item 蓝海枢纽采用漂浮式风电为主体电源，以构网型电池承担快速稳定支撑，以制氢储氢和柔性
        算力参与小时级能量调节，并设置电力、氢能、算力和海洋综合用能等输出端口。该架构能够
        在统一公共母线和计量体系下描述多类设备的功率竞争、状态递推和价值核算。
  \item 第四章建立了包含多源出力、电池能量状态、电解制氢与储氢、绿色算力任务队列、动态PUE、
        公共母线平衡及多目标评价的逐时模型。物理电量、算力服务量和IT耗电采用独立单位计量，
        从而避免在能源与算力核算中混用MWh、MWh-CS和MWh-IT。
  \item 合成48~h算例通过20项自动检查，功率平衡、电池状态、氢库存和服务量核算均满足数值容差。
        在正常典型48~h窗口中，在线协同策略的可再生能源利用率为93.933\%，较纯电资产基准提高
        1.677个百分点，运行净收益提高7.91\%。该差异同时受到资产组合和控制策略影响，尚不能
        单独解释为算法增益。
  \item 在现有价格、效率和最低采购量设定下，制氢路径在正常48~h和年度算例中均未进入最优解。
        该结果说明氢能路径的启用条件对承购价格、设备效率、储运成本和最低采购量较为敏感，后续
        需要结合真实合同和物流方案开展阈值分析。
  \item 年度在线策略完成8,760~h连续求解，可再生能源利用率为64.282\%，关键负荷服务率为
        96.184\%。年度未供能量主要发生在正常资源不足时段，表明当前资源输入、关键负荷与
        160~MWh电池容量组合尚未满足工程可靠性要求。台风过境24~h零出力条件下，电池可供电
        119.7~MWh，约可覆盖6.23个等效关键负荷小时，长时保供仍需其他储能或负荷分级措施。
  \item 现阶段经济结果反映运行层收入与变动成本，环境结果为基于统一排放因子的情景比较。由于
        尚未纳入完整CAPEX、固定运维、设备更换、融资和生命周期清单，相关结果不构成项目投资
        可行性或认证减排结论。
\end{enumerate}

\section{研究贡献与应用价值}

本研究的主要贡献体现在三个方面。第一，提出适用于深远海弱网环境的源储算用系统架构，
明确多源发电、构网型储能、制氢储氢、绿色算力和多类输出端口之间的功能关系。第二，建立
兼顾物理功率、设备状态、任务服务和经济评价的统一调度模型，并在非前视信息条件下开展年度
连续仿真。第三，将技术成熟度、仿真指标、合同条件和投资阶段相联系，形成由岸基联调、海上
示范到预商业化扩展的实施思路。

对于国家能源集团，该研究可用于支持深远海综合能源项目的概念设计、场址与产品组合比选、
技术验证任务分解以及合作伙伴接口讨论。模型保留了设备曲线、资源序列、客户任务、价格和
故障状态的替换接口，可在获得企业数据后用于更新容量方案和运行策略。

\section{研究局限与后续工作}

当前研究仍存在以下局限。资源侧采用公开网格数据和内部潮流序列，尚未完成目标场址的测风、
测光、海温和流速校正；设备侧缺少拟选风机、海缆、PCS、电解槽和液冷系统的完整性能曲线；
算力侧采用聚合任务与标准化服务量，尚未覆盖作业级连续训练、真实光纤流量和客户SLA；经济与
环境评价缺少工程量清单、设备报价、融资条件和完整生命周期数据；小时级调度也不能替代构网
动态、保护配合、海工生存和氢安全分析。

后续研究可按以下顺序推进：首先建立目标海域的多源资源数据库和完整关键负荷清单，并以OEM
设备曲线更新模型；其次开展BESS与长时储能容量敏感性、多小时滚动预测控制和算力任务级调度；
再次通过REGFM、EMT和HIL验证构网稳定、故障穿越与黑启动过程；最后结合潜在电力用户、氢能
承购方和算力客户的合同数据，完成工程投资、现金流和生命周期环境评价。上述工作将推动模型
由方案研究工具逐步转化为示范工程的设计与决策支持工具。
